% preamble-exam-zh.tex — 极简讲解页共享 preamble
% 目标：复刻「题目独立、提示轻框、路线箭头、主体双栏、答案轻收束」的讲义风格。

\documentclass{exam-zh}

% ---------- 基础配置 ----------
\examsetup{
  page/size = a4paper,
  page/show-head = false,
  page/show-foot = false,
  sealline/show = false,
  font = none,
  math-font = none,
}

% ---------- 包 ----------
\usepackage{geometry}
\geometry{top=10mm,bottom=14mm,left=12mm,right=10mm}
\AtBeginDocument{\newgeometry{top=10mm,bottom=14mm,left=12mm,right=10mm}}

\usepackage{xcolor}
\usepackage{needspace}
\usepackage{enumitem}
\usepackage{array}
\usepackage{tabularx}
\usepackage{tcolorbox}
\tcbuselibrary{skins,breakable}
\usepackage{paracol}
\usepackage{graphicx}

% ---------- 打印/屏幕颜色切换 ----------
% 用法：\documentclass[monochrome]{exam-zh} 可降级为灰度。
% 注意：不要使用 \if@xxx 放在 makeatother 外部；这里改成普通条件名。
\newif\ifeduprintcolor
\eduprintcolortrue
\makeatletter
\@ifclasswith{exam-zh}{monochrome}{\eduprintcolorfalse}{}
\makeatother

% ---------- 语义颜色 ----------
\ifeduprintcolor
  \definecolor{edu-blue}{RGB}{31,62,126}
  \definecolor{edu-blue-soft}{RGB}{232,238,250}
  \definecolor{edu-blue-bg}{RGB}{248,250,254}
  \definecolor{edu-ink}{RGB}{30,35,45}
  \definecolor{edu-muted}{RGB}{118,124,136}
  \definecolor{edu-line}{RGB}{209,218,235}
  \definecolor{edu-soft-bg}{RGB}{250,251,253}
  \definecolor{edu-red}{RGB}{168,58,58}
  \definecolor{edu-red-bg}{RGB}{255,247,245}
\else
  \definecolor{edu-blue}{gray}{0.15}
  \definecolor{edu-blue-soft}{gray}{0.90}
  \definecolor{edu-blue-bg}{gray}{0.98}
  \definecolor{edu-ink}{gray}{0.10}
  \definecolor{edu-muted}{gray}{0.45}
  \definecolor{edu-line}{gray}{0.82}
  \definecolor{edu-soft-bg}{gray}{0.97}
  \definecolor{edu-red}{gray}{0.28}
  \definecolor{edu-red-bg}{gray}{0.97}
\fi
\colorlet{edu-diagram-muted}{edu-muted}

% ---------- 全局排版 ----------
\setlength{\parindent}{0pt}
\setlength{\parskip}{0pt}
\linespread{1.16}
\renewcommand{\arraystretch}{1.15}

% ctex 通常提供 \kaishu；这里用它做教师旁白感。
\newcommand{\edukai}{\kaishu}

% ---------- 顶部元信息 ----------
\newcommand{\eduTopBar}[2]{%
  \noindent
  \begin{tabularx}{\linewidth}{@{}Xr@{}}
    {\color{edu-blue}\bfseries #1} &
    {\color{edu-blue}\bfseries #2}
  \end{tabularx}
  \vspace{8mm}
}

% ---------- 轻标题体系 ----------
% 只有真正的大结构才用它；不要给提示/路线滥用标题。
\newcommand{\eduSectionTitle}[1]{%
  \needspace{5\baselineskip}%
  \vspace{2mm}
  {\color{edu-blue}\bfseries\large #1}\par
  \vspace{3mm}
}

\newcommand{\eduSubTitle}[1]{%
  \needspace{4\baselineskip}%
  \vspace{1mm}
  {\color{edu-blue}\bfseries #1}\par
  \vspace{1mm}
}

% ---------- 小问题干：复用 exam (1)(2)(3) 格式 ----------
\newenvironment{eduSubQuestionItem}[1]{%
  \needspace{4\baselineskip}%
  \vspace{2mm}
  \par\noindent
  \hangindent=8mm\hangafter=1
  {\color{edu-blue}\bfseries #1}\enspace
  \begingroup
  \color{edu-ink}
}{%
  \endgroup
  \par
  \vspace{3mm}
}

% ---------- 辅助信息条（解题流程/关键想法/思考）楷体 ----------
\newenvironment{eduAuxItem}[1]{%
  \needspace{4\baselineskip}%
  \vspace{2mm}
  \par\noindent
  \hangindent=8mm\hangafter=1
  {\color{edu-blue}\edukai $\triangleright$~#1}\enspace
  \begingroup
  \color{edu-ink}
  \edukai
}{%
  \endgroup
  \par
  \vspace{4mm}
}

\newcommand{\eduColumnTitle}[2]{%
  {\color{edu-blue}\bfseries #1\hspace{1.5mm}#2}\par
  \vspace{2.5mm}
}

% ---------- 图标：不用外部 icon 字体，保证可移植 ----------
\newcommand{\eduIconBulb}{\(\circ\)}
\newcommand{\eduIconBook}{\(\square\)}
\newcommand{\eduIconCheck}{\(\checkmark\)}
\newcommand{\eduIconPin}{\(\diamond\)}
\newcommand{\eduIconNote}{\(\triangleright\)}

% ---------- 题目区 ----------
% 题目区是页面入口：弱标签 + 一条长蓝线 + 大题干。
\newenvironment{eduProblemBlock}[1][题目]{%
  \needspace{8\baselineskip}
  {\small\color{edu-muted}#1}\par
  \vspace{1.5mm}
  {\color{edu-blue}\hrule height 0.55pt}
  \vspace{4.5mm}
  \begingroup
  \color{edu-ink}
  \large
  \setlength{\parskip}{2.2mm}
  \linespread{1.20}\selectfont
}{%
  \par
  \endgroup
  \vspace{6mm}
}

\newenvironment{eduSubQuestions}{%
  \begin{enumerate}[
    label=(\arabic*),
    leftmargin=8mm,
    itemsep=2.0mm,
    topsep=2.0mm,
    parsep=0pt
  ]
}{%
  \end{enumerate}
}

% ---------- 读题提示：轻框，不做同级标题 ----------
\newtcolorbox{eduReadTipBox}{
  enhanced,
  breakable,
  colback=white,
  colframe=edu-blue!45,
  boxrule=0.45pt,
  arc=1.6mm,
  left=10mm,
  right=4mm,
  top=5mm,
  bottom=5mm,
  before skip=1mm,
  after skip=7mm,
  fontupper=\edukai\normalsize,
  colupper=edu-ink,
  overlay={
    \node[anchor=north west, text=edu-blue] at ([xshift=3.2mm,yshift=-3.2mm]frame.north west)
    {\small\eduIconNote};
  }
}

\newenvironment{eduTipList}{%
  \begin{itemize}[
    label=\textbullet,
    leftmargin=5mm,
    itemsep=1.2mm,
    topsep=0pt,
    parsep=0pt
  ]
}{%
  \end{itemize}
}

% ---------- 解题路线：虚线框 + 文字箭头 ----------
\newenvironment{eduRouteFlow}{%
  \needspace{4\baselineskip}%
  \vspace{1mm}
  \begin{center}
  \normalsize
}{%
  \end{center}
  \vspace{7mm}
}
\newcommand{\eduRouteStep}[1]{%
  {\color{edu-ink}#1}%
}
\newcommand{\eduRouteArrow}{%
  \;$\boldsymbol{\to}$\;%
}

% ---------- 主体讲解：使用 paracol 实现跨页自适应双栏 ----------
\newenvironment{eduDualExplain}{%
  \needspace{6\baselineskip}% 防止标题孤行
  \columnratio{0.26, 0.74}%
  \setlength{\columnsep}{8mm}%
  \columnseprule=0.45pt%
  \colseprulecolor{edu-blue}%
  \begin{paracol}{2}% 开启双栏
}{%
  \end{paracol}% 结束双栏
  \vspace{5mm}
}

\newenvironment{eduExplainLeft}{%
  \setlength{\linewidth}{\columnwidth}%
  \hsize=\columnwidth
  \normalsize\color{edu-ink}%
}{%
  \switchcolumn
}

\newcommand{\eduColumnDivider}{}

\newenvironment{eduExplainRight}{%
  \setlength{\linewidth}{\columnwidth}%
  \hsize=\columnwidth
  \normalsize\color{edu-ink}%
}{%
}

\newtcolorbox{eduSideHintBox}[1]{
  enhanced,
  breakable,
  colback=edu-blue-bg,
  colframe=edu-line,
  boxrule=0.35pt,
  borderline west={1.5pt}{0pt}{edu-blue},
  arc=0.8mm,
  left=2.4mm,
  right=2.4mm,
  top=2mm,
  bottom=2mm,
  before skip=1mm,
  after skip=1.5mm,
  title={\color{edu-blue}\bfseries\small #1},
  coltitle=edu-blue,
  attach title to upper={\par\vspace{0.7mm}},
  fontupper=\small
}

\newtcolorbox{eduSideMistakeBox}[1]{
  enhanced,
  breakable,
  colback=edu-red-bg,
  colframe=edu-red!35,
  boxrule=0.35pt,
  borderline west={1.5pt}{0pt}{edu-red},
  arc=0.8mm,
  left=2.4mm,
  right=2.4mm,
  top=2mm,
  bottom=2mm,
  before skip=1mm,
  after skip=1.5mm,
  title={\color{edu-red}\bfseries\small #1},
  coltitle=edu-red,
  attach title to upper={\par\vspace{0.7mm}},
  fontupper=\small
}

\newtcolorbox{eduSideNoteBox}[1]{
  enhanced,
  breakable,
  colback=edu-soft-bg,
  colframe=edu-line,
  boxrule=0.35pt,
  borderline west={1.5pt}{0pt}{edu-muted},
  arc=0.8mm,
  left=2.4mm,
  right=2.4mm,
  top=2mm,
  bottom=2mm,
  before skip=1mm,
  after skip=1.5mm,
  title={\color{edu-muted}\bfseries\small #1},
  coltitle=edu-muted,
  attach title to upper={\par\vspace{0.7mm}},
  fontupper=\small
}

\newcommand{\eduSideHint}[2]{%
  \begin{eduSideHintBox}{#1}#2\end{eduSideHintBox}%
}
\newcommand{\eduSideMistake}[2]{%
  \begin{eduSideMistakeBox}{#1}#2\end{eduSideMistakeBox}%
}
\newcommand{\eduSideNote}[2]{%
  \begin{eduSideNoteBox}{#1}#2\end{eduSideNoteBox}%
}

\newcommand{\eduExplainStep}[3]{%
  \needspace{3\baselineskip}%
  \vspace{1.2mm}%
  {\color{edu-blue}\bfseries step#1.\enspace #2}\par
  \vspace{0.8mm}%
  \begingroup
  \color{edu-ink}#3\par
  \endgroup
  \vspace{2.2mm}%
}

\newenvironment{eduNumberList}{%
  \begin{enumerate}[
    label=\arabic*.,
    leftmargin=6mm,
    itemsep=4.8mm,
    topsep=0pt,
    parsep=0pt
  ]
}{%
  \end{enumerate}
}

\newenvironment{eduBulletList}{%
  \begin{itemize}[
    label=\textbullet,
    leftmargin=5mm,
    itemsep=1.2mm,
    topsep=0pt,
    parsep=0pt
  ]
}{%
  \end{itemize}
}

% ---------- 底部双栏：答案 + 方法提醒 ----------
\newenvironment{eduBottomGrid}{%
  \columnratio{0.32, 0.68}%
  \setlength{\columnsep}{11mm}%
  \columnseprule=0pt%
  \begin{paracol}{2}%
}{%
  \end{paracol}%
}

\newenvironment{eduSummaryLeft}{%
}{%
}

\newcommand{\eduSummaryDivider}{%
  \switchcolumn
}

\newenvironment{eduSummaryRight}{%
}{%
}

% ---------- 答案/方法提醒：轻量收束 ----------
\newtcolorbox{eduSummaryBlock}[2]{
  enhanced,
  breakable,
  colback=edu-blue-bg,
  colframe=edu-line,
  boxrule=0.45pt,
  arc=1.2mm,
  left=3.5mm,
  right=3.5mm,
  top=3mm,
  bottom=3mm,
  before skip=1mm,
  after skip=2.5mm,
  title={\color{edu-blue}\bfseries #1\hspace{1.5mm}#2},
  coltitle=edu-blue,
  attach title to upper={\par\vspace{1.5mm}},
  fontupper=\normalsize
}

% ---------- 例题单栏解答卡片 ----------
\newtcolorbox{eduSolutionBox}[1]{
  enhanced,
  breakable,
  colback=white,
  colframe=edu-blue!40,
  boxrule=0.55pt,
  arc=1.2mm,
  left=4mm,
  right=4mm,
  top=3.5mm,
  bottom=3.5mm,
  before skip=2mm,
  after skip=3mm,
  title={\color{edu-blue}\bfseries \eduIconBook\hspace{1.5mm}#1},
  coltitle=edu-blue,
  attach title to upper={\par\vspace{1.5mm}},
  fontupper=\normalsize
}

% ---------- 变式训练：完整题目 + 学生留白 ----------
\newtcolorbox{eduVariationTraining}[1]{
  enhanced,
  breakable,
  colback=white,
  colframe=edu-blue!45,
  boxrule=0.5pt,
  borderline west={1.8pt}{0pt}{edu-blue},
  arc=1.2mm,
  left=4mm,
  right=4mm,
  top=3.2mm,
  bottom=3.2mm,
  before skip=2mm,
  after skip=4mm,
  title={\color{edu-blue}\bfseries #1},
  coltitle=edu-blue,
  attach title to upper={\par\vspace{1.8mm}},
  fontupper=\normalsize,
  colupper=edu-ink
}

\newenvironment{eduPlainList}{%
  \begin{enumerate}[
    label=(\arabic*),
    leftmargin=7mm,
    itemsep=1.2mm,
    topsep=0pt,
    parsep=0pt
  ]
}{%
  \end{enumerate}
}

% ---------- 兼容旧 block 的轻量盒子 ----------
\newtcolorbox{eduSoftBox}{
  enhanced,
  breakable,
  colback=edu-soft-bg,
  colframe=edu-line,
  boxrule=0.35pt,
  arc=1mm,
  left=3mm,
  right=3mm,
  top=2mm,
  bottom=2mm,
  before skip=1mm,
  after skip=2mm,
  fontupper=\small
}

\newtcolorbox{eduMistakeBox}{
  enhanced,
  breakable,
  colback=edu-red-bg,
  colframe=edu-red!40,
  boxrule=0.35pt,
  arc=1mm,
  left=3mm,
  right=3mm,
  top=2.5mm,
  bottom=2.5mm,
  before skip=1mm,
  after skip=2mm,
  fontupper=\normalsize
}

\newtcolorbox{eduLegacySolution}{
  enhanced,
  breakable,
  colback=edu-soft-bg,
  colframe=edu-line,
  boxrule=0.35pt,
  arc=1mm,
  left=3mm,
  right=3mm,
  top=2mm,
  bottom=2mm,
  before skip=-2mm,
  after skip=4mm,
  fontupper=\small
}

\newtcolorbox{eduTeacherNote}{
  enhanced,
  breakable,
  colback=edu-soft-bg,
  colframe=edu-line,
  boxrule=0pt,
  borderline west={1.5pt}{0pt}{edu-muted},
  left=2.5mm,
  right=2.5mm,
  top=2mm,
  bottom=2mm,
  before skip=1mm,
  after skip=2mm,
  fontupper=\footnotesize,
  colupper=edu-muted
}

% ---------- 答题区 ----------
\newcommand{\answerarea}[1][40mm]{%
  \vspace{3mm}%
  \begin{tcolorbox}[
    enhanced,
    breakable,
    colback=white,
    colframe=edu-line,
    boxrule=0.55pt,
    arc=0.8mm,
    height=#1,
    valign=top,
  ]
  \end{tcolorbox}%
}

\newcommand{\diagramcolfigure}[3]{%
  \begin{minipage}[t]{#2}
    \vspace{0pt}\centering
    \includegraphics[width=\linewidth]{\detokenize{#1}}
    \if\relax\detokenize{#3}\relax\else
      \par\vspace{1mm}{\scriptsize\color{edu-diagram-muted}#3}
    \fi
  \end{minipage}%
}

\newcommand{\diagramrowitem}[4]{%
  \begin{minipage}[t]{#2}
    \vspace{0pt}\centering
    \includegraphics[width=\linewidth]{\detokenize{#1}}
    \if\relax\detokenize{#3}\relax\else
      \par\vspace{1mm}{\scriptsize\bfseries #3}
    \fi
    \if\relax\detokenize{#4}\relax\else
      \par\vspace{0.5mm}{\scriptsize\color{edu-diagram-muted}#4}
    \fi
  \end{minipage}%
}

\newcommand{\answerareawithdiagramcol}[4]{%
  \vspace{3mm}%
  \begin{tcolorbox}[
    enhanced,
    breakable,
    colback=white,
    colframe=edu-line,
    boxrule=0.55pt,
    arc=0.8mm,
    height=#1,
    valign=top,
  ]
  \begin{minipage}[t]{\dimexpr\linewidth-#3-6mm\relax}
    \vspace{0pt}
  \end{minipage}\hfill
  \diagramcolfigure{#2}{#3}{#4}
  \end{tcolorbox}%
}

\newcommand{\answerareawithdiagram}[4]{%
  \answerareawithdiagramcol{#1}{#2}{#3}{#4}%
}

\examsetup{
  question/show-answer = true,
  fillin/show-answer = true,
  solution/show-solution = show-stay,
}

\begin{document}

\eduSectionTitle{对称轴反求：y 相等的两点关于对称轴对称}


\begin{eduSolutionBox}{核心结论：y 相等的两点关于对称轴对称}
  \begingroup
  \setlength{\parskip}{2.5mm}
\textbf{定理：} 若二次函数 $y=f(x)$ 上两点 $(x_1,y_0),(x_2,y_0)$ 满足 $f(x_1)=f(x_2)$ 且 $x_1\ne x_2$，则对称轴为 $x=h=\dfrac{x_1+x_2}{2}$。\par
\textbf{记忆：} 对称轴 = 两个等 $y$ 点横坐标的\textbf{平均数}（不是差，不是半差）。\par
\textbf{结合一般式：} $y=ax^2+bx+c$ 的对称轴 $h=-\dfrac{b}{2a}$，联立即得 $b=-2ah$。\par
\begin{tabular}{p{0.34\linewidth}p{0.56\linewidth}}
题面信号 & 入口动作\\
$f(m)=f(n)$ 或 $(m,y_0),(n,y_0)$ & 取 $h=\dfrac{m+n}{2}$\\
四个点含两对等 $y$ & 先分组，再用上式分别列方程\\
已知对称轴 $h$，求 $b$ 与 $a$ 关系 & 用 $b=-2ah$\\
\end{tabular}
\par
  \endgroup
\end{eduSolutionBox}




\begin{eduProblemBlock}[例题 1（两等点求对称轴）]
已知二次函数 $y=f(x)$ 满足 $f(1)=f(5)$，求其图象的对称轴。

\end{eduProblemBlock}





\begin{eduAuxItem}{思路导航}
识别 $y$ 相等的两点$\;\to\;$取横坐标平均数得对称轴\end{eduAuxItem}




\begin{eduSubQuestionItem}{例题 1}
求二次函数 $y=f(x)$（$f(1)=f(5)$）的对称轴。\end{eduSubQuestionItem}

\begin{eduDualExplain}
  \begin{eduExplainLeft}
    \eduColumnTitle{\eduIconBulb}{提示与易错}
        \eduSideHint{为什么用平均数？}{对称轴两侧等距的点函数值相等，所以对称轴正好在两点正中间，即平均数位置。}
        \eduSideMistake{别算成差}{对称轴是 $\dfrac{1+5}{2}=3$，不是 $5-1=4$，也不是 $\dfrac{5-1}{2}=2$。}
  \end{eduExplainLeft}

  \begin{eduExplainRight}
    \eduColumnTitle{\eduIconBook}{解答}
      \eduExplainStep{1}{识别 $y$ 相等的两点}{$f(1)=f(5)$ 表示 $(1,y_0)$ 与 $(5,y_0)$ 纵坐标相等。}
      \eduExplainStep{2}{取横坐标平均数得对称轴}{$h=\dfrac{1+5}{2}=3$。}
  \end{eduExplainRight}
\end{eduDualExplain}




\begin{eduProblemBlock}[例题 2（必考四点变式）]
抛物线上有四个点 $(1,1),(3,1),(4,-1),(a,-1)$，求 $a$ 的值。

\end{eduProblemBlock}





\begin{eduAuxItem}{思路导航}
把 $y$ 相等的点分组$\;\to\;$用第一组定对称轴$\;\to\;$第二组关于对称轴对称，列方程求 $a$\end{eduAuxItem}




\begin{eduSubQuestionItem}{例题 2}
抛物线上四点 $(1,1),(3,1),(4,-1),(a,-1)$，求 $a$。\end{eduSubQuestionItem}

\begin{eduDualExplain}
  \begin{eduExplainLeft}
    \eduColumnTitle{\eduIconBulb}{提示与易错}
        \eduSideHint{先分组再动手}{四个点看起来乱，关键是先按 $y$ 值分成两组：$y=1$ 一组，$y=-1$ 一组。}
        \eduSideHint{两组共用同一条对称轴}{同一条抛物线只有一条对称轴，所以两组等 $y$ 点都关于它对称。}
        \eduSideMistake{别四个点混在一起硬解}{不要把四个点全部代入 $y=ax^2+bx+c$ 列方程组，计算爆炸且极易算错。}
  \end{eduExplainLeft}

  \begin{eduExplainRight}
    \eduColumnTitle{\eduIconBook}{解答}
      \eduExplainStep{1}{把 $y$ 相等的点分组}{$(1,1)$ 与 $(3,1)$ 一组（$y=1$）；$(4,-1)$ 与 $(a,-1)$ 一组（$y=-1$）。}
      \eduExplainStep{2}{用第一组定对称轴}{由 $(1,1),(3,1)$ 得 $h=\dfrac{1+3}{2}=2$。}
      \eduExplainStep{3}{第二组关于对称轴对称，列方程求 $a$}{$(4,-1)$ 与 $(a,-1)$ 关于 $x=2$ 对称，故 $\dfrac{4+a}{2}=2$，解得 $a=0$。}
    \vspace{1mm}
    {\color{edu-blue}\bfseries\small 一句话收束}\par\vspace{1mm}
    \begin{eduBulletList}
      \item 见等 $y$ 即取平均数；含未知坐标的那组就用来反解未知数。    \end{eduBulletList}
  \end{eduExplainRight}
\end{eduDualExplain}




\begin{eduProblemBlock}[例题 3（结合一般式求 $b$）]
已知二次函数 $y=x^2+bx+c$ 满足 $f(1)=f(5)$。
\begin{enumerate}[label=(\arabic*)]
  \item 求该函数图象的对称轴；
  \item 求 $b$ 的值。
\end{enumerate}

\end{eduProblemBlock}





\begin{eduAuxItem}{思路导航}
(1) 取平均数求对称轴$\;\to\;$(2) 用 $h=-\dfrac{b}{2a}$ 求 $b$\end{eduAuxItem}




\begin{eduSubQuestionItem}{(1)}
求该函数图象的对称轴；\end{eduSubQuestionItem}

\begin{eduDualExplain}
  \begin{eduExplainLeft}
    \eduColumnTitle{\eduIconBulb}{提示与易错}
        \eduSideHint{看到 $f(m)=f(n)$ 想什么}{想对称性：直接取平均数 $\dfrac{m+n}{2}$，无需先求 $b,c$。}
  \end{eduExplainLeft}

  \begin{eduExplainRight}
    \eduColumnTitle{\eduIconBook}{解答}
      \eduExplainStep{1}{(1) 取平均数求对称轴}{由 $f(1)=f(5)$ 得 $h=\dfrac{1+5}{2}=3$。}
  \end{eduExplainRight}
\end{eduDualExplain}




\begin{eduSubQuestionItem}{(2)}
求 $b$ 的值。\end{eduSubQuestionItem}

\begin{eduDualExplain}
  \begin{eduExplainLeft}
    \eduColumnTitle{\eduIconBulb}{提示与易错}
        \eduSideHint{对称轴与 $a,b$ 的桥梁}{一般式对称轴 $h=-\dfrac{b}{2a}$ 是连接 $h$ 与 $a,b$ 的唯一公式；已知 $h$ 和 $a$ 即可解 $b$。}
        \eduSideMistake{别漏负号}{$h=-\dfrac{b}{2a}$ 中等号右边自带负号，代入时 $-\dfrac{b}{2}=3$，不要写成 $\dfrac{b}{2}=3$。}
  \end{eduExplainLeft}

  \begin{eduExplainRight}
    \eduColumnTitle{\eduIconBook}{解答}
      \eduExplainStep{1}{(2) 用 $h=-\dfrac{b}{2a}$ 求 $b$}{$a=1$，由 $-\dfrac{b}{2a}=3$ 得 $-\dfrac{b}{2}=3$，解得 $b=-6$。}
  \end{eduExplainRight}
\end{eduDualExplain}




\begin{eduMistakeBox}
\eduSubTitle{易错提醒}\begin{itemize}
  \item 忽略“$y$ 相等的两点关于对称轴对称”这条对称性，转去硬求解析式。
  \item 把对称轴 $h=\dfrac{x_1+x_2}{2}$ 算成差 $x_2-x_1$ 或半差 $\dfrac{x_2-x_1}{2}$。
  \item 四点题没有先按 $y$ 值分组，把四个点搅在一起处理。
  \item 用 $h=-\dfrac{b}{2a}$ 时漏掉等号右边的负号。
\end{itemize}
\end{eduMistakeBox}




\begin{eduSummaryBlock}{\eduIconPin}{方法提醒}
\begin{eduBulletList}
  \item 见 $f(m)=f(n)$ 或两点纵坐标相等 $\Rightarrow$ 对称轴 $h=\dfrac{m+n}{2}$。  \item 四点题先按 $y$ 值分组，每组各用一次平均数。  \item 结合一般式求 $b$：用 $b=-2ah$（其中 $h$ 由对称性得到）。\end{eduBulletList}
\end{eduSummaryBlock}




\begin{eduBottomGrid}
  \begin{eduSummaryLeft}
    \begin{eduSummaryBlock}{\eduIconCheck}{答案}
    \begin{eduPlainList}
      \item 例 1：对称轴 $x=3$      \item 例 2：$a=0$      \item 例 3：(1) 对称轴 $x=3$；(2) $b=-6$    \end{eduPlainList}
    \end{eduSummaryBlock}
  \end{eduSummaryLeft}
  \eduSummaryDivider
  \begin{eduSummaryRight}
    \begin{eduSummaryBlock}{\eduIconPin}{一句话}
    \begin{eduBulletList}
      \item $y$ 相等的两点关于对称轴对称 $\Rightarrow$ 对称轴 $=$ 横坐标平均数。    \end{eduBulletList}
    \end{eduSummaryBlock}
  \end{eduSummaryRight}
\end{eduBottomGrid}




\end{document}
